TESS (Transiting Exoplanet Survey Satellite) is NASA's most recent exoplanet
hunter mission. It has been surveying the southern sky to locate exoplanets
around the brightest and closest stars, along with a large number of
time-variable phenomena (pulsating and eclipsing variable stars, supernovae,
stellar flares, and asteroids among others).

On the separate sheet you will find graphs with light curve plots of 8
periodic variable stars (intrinsic, eclipsing and rotational; namely FO Eri,
RW Dor, 24 Eri, TIC 147272181, ST Pic, UY Eri, VV Ori, AH Col) from the TESS
target list, numbered from 1 to 8. The horizontal axis is
$\mathrm{BJD} - 2\,400\,000$ in days, and the vertical axis is $V$ in
magnitudes, where BJD is the Barycentric Julian Date, and $V$ is the visual
brightness.

\begin{description}
\item[(a)] Below is a list of variable star types. Match each light curve of
  the star with its variability type by writing its number into the rectangle
  on the answer sheet. \hfill[8\,p]
  \begin{itemize}
  \item Heart-beat star
  \item RR Lyrae type (RRab subclass) pulsating variable star
  \item Eclipsing binary of Algol type (semi-detached) with a pulsating
    component
  \item $\alpha^2$ CVn pulsating variable star
  \item W Vir type (Population II) Cepheid pulsating variable star
  \item Detached eclipsing binary with strong reflection effect
  \item Contact eclipsing binary of W UMa type
  \item Rotationally variable (spotted) star
  \end{itemize}
\item[(b)] Based on the light curve plots, estimate the periods of each
  variable star in days. Give your answers on the answer sheet up to 2
  decimal places. Periods in the range $\pm5$\% of the true periods are
  acceptable. \hfill[16\,p]
\item[(c)] What type of astronomical object produced the TESS light curve
  shown in the figure below? Mark the letter of your answer with $\times$ on
  the answer sheet. \hfill[1\,p]
  \begin{description}
  \item[(A)] Microlensing event
  \item[(B)] Saturated galaxy due to the proximity of Mars
  \item[(C)] Comet's coma close to the edge of the camera field
  \item[(D)] Supernova in a distant galaxy
  \item[(E)] Superflare on a supergiant star
  \end{description}
\end{description}
